\documentclass[11pt,a4paper]{article}
\usepackage[T1]{fontenc}
\usepackage{lmodern}
\usepackage[margin=25mm]{geometry}
\usepackage{amsmath,amssymb,amsthm}
\usepackage{booktabs}
\usepackage{xcolor}
\usepackage[colorlinks=true,linkcolor=blue!50!black,urlcolor=blue!50!black]{hyperref}

\newtheorem{proposition}{Proposition}[section]
\newtheorem{lemma}[proposition]{Lemma}
\theoremstyle{remark}
\newtheorem*{remark}{Remark}

\setlength{\emergencystretch}{3em}
\newcommand{\code}[1]{\texttt{\detokenize{#1}}}
\newcommand{\takeaway}[1]{\par\medskip\noindent\fbox{\parbox{\dimexpr\linewidth-2\fboxsep-2\fboxrule\relax}{\textbf{Plain-English takeaway.} #1}}\par\medskip}
\newcommand{\E}{\mathbb{E}}
\newcommand{\Var}{\operatorname{Var}}
\newcommand{\Cov}{\operatorname{Cov}}
\newcommand{\plim}{\operatorname*{plim}}
\newcommand{\logit}{\operatorname{logit}}
\newcommand{\rs}{r^{*}}
\newcommand{\one}{\mathbf{1}}
\DeclareMathOperator*{\argmin}{arg\,min}

\title{Mortgage price-elasticity model:\\ the mathematics behind the code}
\author{Technical companion to \texttt{docs/overview.md}}
\date{}

\begin{document}
\maketitle

\section*{How to read this document}
Section~$N$ expands step~$N$ of \texttt{docs/overview.md}. Every formula names the function that
implements it, as \code{module.function} under \code{src/mortgage_elasticity/}. Defaults quoted come from
\code{config.py} (\code{Config}). The text describes what the code does now. Where
the code makes an approximation or a judgment call, the text says so. Each section ends with a boxed
takeaway that you can pass on to business users.

\paragraph{What changed from the earlier design.} The earlier pipeline produced five sensitivity
estimates (two-way fixed effects, event studies, a per-segment posterior, a pooled posterior and an
optional hierarchical model) and a separate Kalman step for natural share. It is replaced by one
Bayesian state-space model of lodgement counts (\S\S\ref{sec:alpha}--\ref{sec:fit}). Two-way fixed
effects and event studies survive only as diagnostics (\S\ref{sec:check}).

\tableofcontents

%=====================================================================
\section{Setup and notation}\label{sec:setup}

\paragraph{Data.} Segments $k=1,\dots,K$ (up to 20) are purpose (OO, INV) $\times$ repayment (PI, IO)
$\times$ LVR band ($\le$60, 60--70, 70--80, 80--90, $>$90), defined identically in our data and the
market data and named \code{PURPOSE_REPAYMENT_BAND} (\code{panel.segment_dims}). Periods
$t=1,\dots,n$ are weeks, labelled by their \emph{end} date, broker channel only. The panel is
balanced, which \code{panel.load_panel} enforces. For each $(k,t)$ we observe market lodgements
$M_{kt}$ (\code{lodg_mkt}), our lodgements $L_{kt}$ (\code{lodg_us}), our applications $A_{kt}$
(\code{apps}), and two cumulative price paths since the panel start in basis points (bp), with $+$
meaning dearer: $H_{kt}$, the carded rate (\code{headline_cum_bps}), and $D_{kt}$, the average
discretionary discount (\code{discretionary_cum_bps}; $+$ means less discount; treated as zero, and
logged, if the column is absent).

\paragraph{Counts, not ratios.} The observed share is $s_{kt}=L_{kt}/M_{kt}$, but the model is
written for the count $L_{kt}$ given $M_{kt}$ (\S\ref{sec:fit}). The reason is the sampling noise.
If $L\mid M\sim\mathrm{Bin}(M,\pi)$ then $\Var(s)=\pi(1-\pi)/M$ and, by the delta method with
$\tfrac{d}{ds}\logit s=1/(s(1-s))$,
\begin{equation}\label{eq:sampling}
  \Var(\logit s_{kt})\approx\frac{1}{M_{kt}\,\pi_{kt}(1-\pi_{kt})}.
\end{equation}
For $M=4000$, $\pi=0.10$ the standard deviation is $0.053$; for a thin segment with $M=60$ it is
$0.43$, and $L=0$ occurs in $0.9^{60}\approx0.2\%$ of weeks at $\pi=0.1$ and in $5\%$ of weeks at
$\pi=0.05$. A regression on $\logit s$ must drop or patch those weeks and needs weights it does not
know. A count likelihood handles both: a week with $L=0$ is an ordinary observation, and the noise
scales with $M$ automatically. \code{panel.load_panel} therefore rejects $L>M$, negative or
fractional counts and duplicate rows, but keeps $L=0$. Weeks with $M=0$ carry no information and are
left out of the likelihood (\code{Prepared.obs}) and counted in the run log.

\paragraph{Conversion assumption.} Business users plan in applications, but the model is
estimated on lodgements. The bridge is $A_{kt}=c_k L_{kt}$ with a conversion ratio $c_k$ that is
constant over the forecast horizon. \code{panel.conversion_report} estimates it as trailing 13-week
$\sum A/\sum L$ and checks the assumption by reporting the coefficient of variation of
$A_{kt}/L_{kt}$ and its OLS slope on $t$ divided by its mean
(\code{conv_trend_per_period_pct}). A material trend means that a change in lodgement share is not
a change in application share.

\paragraph{Event log.} Each logged change has a date, the segments affected (or \code{ALL}), and a
type: \code{price} for an own price move, \code{other} for anything else. A date inside a week
lands in that week. Price events add nothing to the model, because their effect is carried by $x$;
they are used by the diagnostics and backtests. Other events get a level jump (\S\ref{sec:alpha}).

\paragraph{Notation.} Symbols are used with one meaning throughout.

\begin{center}
\footnotesize
\begin{tabular}{@{}lll@{}}
\toprule
Symbol & Meaning & Code \\
\midrule
$k$, $t$, $K$, $n$ & segment, week, number of segments, number of weeks & \code{segment}, \code{period}\\
$M_{kt}$, $L_{kt}$, $s_{kt}$ & market and own lodgements; observed share $L/M$ & \code{lodg_mkt}, \code{lodg_us}\\
$\pi_{kt}$ & model share (choice probability) & \code{Fit.pi}\\
$H_{kt}$, $D_{kt}$ & our cumulative carded / discretionary change (bp) & \code{headline_cum_bps}\\
$q_{jkt}$ & competitor $j$'s effective rate, $100\cdot$rate$-$adj (bp) & \code{eff_bps}\\
$\rs_{kt}$ & soft-min market benchmark (bp) & \code{r_star_bps}\\
$T_k$ & soft-min temperature (bp), solved so that $T_k=1/\beta_k$ & \code{T_bps}\\
$w_{jkt}$ & lender $j$'s weight in $\rs$ (softmax) & \code{weight}\\
$\kappa$ & discretionary shading (fixed, $0.5$) & \code{kappa}\\
$v_{kt}$ & unlagged price gap (bp, $+$ = dearer) & \code{gap_raw}\\
$h$, $\phi$ & half-life in weeks; weekly adjustment speed $\phi=1-2^{-1/h}$ & \code{h}\\
$W(h)$ & lag kernel, $W_{tj}=\phi(1-\phi)^{t-j}$ for $j\le t$ & \code{model.lag_kernel}\\
$x_{kt}$ & lagged (perceived) price gap & \code{Fit.x}\\
$\alpha_k$ & segment constant of natural share (logit) & \code{alpha}\\
$g_t$ & bank-wide trend, random walk, $g_1=0$ & \code{sigma_g}\\
$\nu_{kt}$ & segment trend, random walk, $\nu_{k1}=0$, innovations zero-sum over $k$ & \code{sigma_v}\\
$J_{kt}$ & sum of level jumps $j_e$ for logged non-price changes & \code{jump}\\
$\ell_{kt}$ & natural-share level $\alpha_k+g_t+\nu_{kt}+J_{kt}$ & \code{level}\\
$\beta_k$ & price sensitivity of the logit share, per bp & \code{beta}\\
$u_k$ & uplift per bp at baseline, $u_k=\beta_k(1-s_{0k})$ & \code{u}\\
$s_{0k}$ & pooled share $\sum_tL_{kt}/\sum_tM_{kt}$ (a fixed constant) & \code{Prepared.s0}\\
$\mu_u$, $a_{\rm inv}$, $a_{\rm io}$, $c^{\rm lvr}_b$ & pooling parameters of $\log u_k$ & \code{mu_u}, \code{a_inv}, \dots\\
$c$ & beta-binomial concentration & \code{c}\\
$\Lambda(z)$ & logistic function $1/(1+e^{-z})$ & \\
\bottomrule
\end{tabular}
\end{center}

The segment trend is written $\nu_{kt}$ here, and \code{seg_trend} in code, to keep it apart from
the unlagged gap $v_{kt}$ (\code{gap_raw}).

\takeaway{We count our lodgements and the market's in each segment every week, and assume
applications move in step with lodgements. The model works with the counts, so a thin week counts
for less than a busy one and a week with no lodgements is still usable.}

%=====================================================================
\section{Customer choice and the market benchmark}\label{sec:choice}

\subsection{From random utility to the multinomial logit}
In segment $k$ and period $t$, borrower $i$ chooses one lender $j$ from the set $\{us\}\cup\mathcal J$,
where $\mathcal J$ is the set of competitors. The borrower gets utility
\[
  U_{ij}=a_j-b\,r_j+\epsilon_{ij},\qquad \epsilon_{ij}\overset{iid}{\sim}\text{Gumbel}(0,1),\quad
  F(\epsilon)=\exp(-e^{-\epsilon}).
\]
Here $a_j$ is lender $j$'s non-price appeal, $r_j$ its rate in bp, and $b>0$ the price
coefficient. Indices $k,t$ are dropped in this subsection.

\begin{proposition}[Multinomial logit]\label{prop:mnl}
With $V_j=a_j-br_j$, $\Pr(i\text{ chooses }j)=e^{V_j}/\sum_{l}e^{V_l}$.
\end{proposition}
\begin{proof}
Lender $j$ is chosen iff $\epsilon_{il}<V_j-V_l+\epsilon_{ij}$ for all $l\ne j$. Condition on
$\epsilon_{ij}=\epsilon$ and use independence:
\[
\Pr(j)=\int_{-\infty}^{\infty}e^{-\epsilon}e^{-e^{-\epsilon}}\prod_{l\ne j}\exp\!\big(-e^{-\epsilon}e^{V_l-V_j}\big)\,d\epsilon
=\int_{-\infty}^{\infty}e^{-\epsilon}\exp\!\big(-e^{-\epsilon}S_j\big)\,d\epsilon,
\]
where $S_j=\sum_l e^{V_l-V_j}$ (the $l=j$ term is the factor $e^{-e^{-\epsilon}}$). Substituting
$z=e^{-\epsilon}$ gives $\int_0^\infty e^{-zS_j}dz=1/S_j$.
\end{proof}

With many borrowers, our share is the choice probability plus the sampling noise of
\eqref{eq:sampling}. The market here is all lodgements in the segment. There is no
``outside good'' (not borrowing), so the market total is taken as given, as in \S\ref{sec:forecast}.

\subsection{Our share against the field}
Write $P_{us}$ for our choice probability. Then
\[
\logit P_{us}=\log\frac{e^{V_{us}}}{\sum_{j\in\mathcal J}e^{V_j}}
= a_{us}-b\,r_{us}-\log\sum_{j\in\mathcal J}e^{a_j-br_j}.
\]
Define the \emph{quality-adjusted} rate $q_j=r_j-a_j/b$ and the temperature $T=1/b$. Then
\begin{equation}\label{eq:binary}
\boxed{\;\logit P_{us}=a_{us}-b\,(r_{us}-\rs),\qquad
\rs=-\frac1b\log\sum_{j\ne us}e^{a_j-br_j}=-T\log\sum_{j\ne us}e^{-q_j/T}.\;}
\end{equation}
This identity is exact. All competitors collapse into one number $\rs$, the logit
\emph{inclusive value} expressed in rate units. Our logit share is linear in our rate's gap to it, with
slope $b$. In the model $b$ is called $\beta_k$, and the two are the same parameter. Our own appeal
$a_{us}$ is everything that is not price: it becomes the natural-share level $\ell_{kt}$ of \S\ref{sec:alpha}.

\subsection{$\rs$ is a soft minimum}
\begin{proposition}\label{prop:softmin}
Let $q_1,\dots,q_N$ be quality-adjusted rates, $q_{\min}=\min_j q_j$, $T>0$ and
$\rs(q)=-T\log\sum_j e^{-q_j/T}$. Then:
\begin{enumerate}
\item[(a)] $q_{\min}-T\log N\le \rs\le q_{\min}$, with the right-hand equality iff $N=1$.
\item[(b)] $\rs\to q_{\min}$ as $T\to0$, uniformly, with error at most $T\log N$.
\item[(c)] Translation equivariance: $\rs(q+c\one)=\rs(q)+c$.
\item[(d)] Adding a lender with rate $q_{N+1}$ strictly lowers $\rs$, by
$T\log\!\big(1+e^{(\rs-q_{N+1})/T}\big)$.
\item[(e)] If all $N$ lenders are tied at $q$, then $\rs=q-T\log N$.
\item[(f)] $\partial\rs/\partial q_j=w_j:=e^{-q_j/T}/\sum_l e^{-q_l/T}$, which are softmax weights summing to one.
\item[(g)] Variational form: $\rs=\min_{w\in\Delta_N}\big\{\sum_j w_jq_j-T\,\mathcal H(w)\big\}$ with
$\mathcal H(w)=-\sum_j w_j\log w_j$. The minimiser is the softmax $w$ of (f).
\end{enumerate}
\end{proposition}
\begin{proof}
(a) $e^{-q_{\min}/T}\le\sum_j e^{-q_j/T}\le Ne^{-q_{\min}/T}$. Apply $-T\log(\cdot)$, which is decreasing.
(b) follows from (a). (c) $\sum_j e^{-(q_j+c)/T}=e^{-c/T}\sum_je^{-q_j/T}$.
(d) The new sum is $\Sigma+e^{-q_{N+1}/T}$ with $\Sigma=e^{-\rs/T}$, so
$\rs_{\rm new}=-T\log\Sigma-T\log(1+e^{-q_{N+1}/T}/\Sigma)$.
(e) $\sum_je^{-q/T}=Ne^{-q/T}$. (f) Differentiate.
(g) For any $w$ in the simplex, concavity of $\log$ (Jensen) gives
$\log\sum_j e^{-q_j/T}=\log\sum_j w_j\frac{e^{-q_j/T}}{w_j}\ge\sum_j w_j\big(-q_j/T-\log w_j\big)$,
with equality iff $e^{-q_j/T}/w_j$ is constant, i.e.\ $w$ is the softmax. Multiply by $-T$.
\end{proof}

\paragraph{Interpretation.} By (g), $\rs$ is a softmax-weighted average rate minus a ``variety bonus''
$T\mathcal H(w)$. With small $T$ (price-sensitive borrowers) it hugs the cheapest lender. With large $T$
the borrower values the spread of options. By (f), a lender's influence on the benchmark is its
softmax weight. A lender priced $3T$ above the cheapest carries $e^{-3}\approx5\%$ of the
cheapest one's weight.

\paragraph{Why entry and exit matter.} By (d) and (e), $\rs$ depends on the \emph{number} of lenders as
well as their rates. At an uplift of 3\%/bp, $T\approx30$\,bp. With 7 tied lenders
the offset is $T\log7\approx58$\,bp, and going from 6 to 8 lenders lowers $\rs$ by
$T\log(8/6)\approx9$\,bp with no rate moving. A constant offset is absorbed by $\alpha_k$
(\S\ref{sec:alpha}). A \emph{change} in $N$ moves $x$. In the model this movement is real: a new
option takes share under IIA. If the change is a data artefact, such as a lender missing from the file or a
forward-fill running out, it is a spurious price move. \code{panel.add_benchmark} warns when $N$
varies within a segment, when $N<$ \code{min_lenders} (3), and when a rate sits more than 100\,bp below the
period median.

\subsection{Why not a weighted average of competitor rates?}
A fixed-weight average $\bar q=\sum_j\pi_j q_j$ contradicts \eqref{eq:binary} in three ways.
(i) Its sensitivity to lender $j$ is a fixed $\pi_j$, whereas the model's is $w_j$, which is close to $0$ for an
uncompetitive lender. An expensive lender's repricing moves $\bar q$ but not our share.
(ii) Adding an expensive lender raises $\bar q$, while the model says it (weakly) lowers $\rs$.
(iii) For equal weights,
$\rs+T\log N=-T\log\big(\tfrac1N\sum_je^{-q_j/T}\big)\to\bar q$ as $T\to\infty$, by a first-order expansion
of the exponential. So the plain average is the benchmark of borrowers who barely react to rate
differences between lenders. That contradicts the premise that $\beta$ is large enough to matter.

\subsection{What the code computes}
\begin{itemize}
\item \textbf{Effective rate} (\code{panel.load_competitors}): $q_{jkt}=100\cdot\code{rate_pct}-\code{adj_bps}$,
with $\code{adj_bps}>0$ meaning value to the customer such as cashback. In \eqref{eq:binary} this sets
$a_j/b=\code{adj}_j$. All other non-price appeal is assumed equal across competitors, and any common part
of it shifts $\rs$ by a constant. \code{adj_bps} is maintained by hand and is the only place where known
competitor discounting can enter; it moves the price variable one for one.
\item \textbf{Eligibility}: rows with \code{eligible}=FALSE get $\mathrm{eff}=+\infty$ and lenders in
\code{exclude_lenders} are dropped. The $+\infty$ marker gives the lender weight $e^{-\infty/T}=0$ in \eqref{eq:binary}, and it
stops \code{panel._wide}, which forward-fills missing rates for up to \code{comp_ffill_periods} (8) periods, from carrying
the last eligible rate into an ineligible period. In a change log, a period's rate averages only the eligible days; a
week that is ineligible on every day is absent.
\item \textbf{Temperature}: the model requires $T_k=1/\beta_k$, so $T_k$ is in bp. Because $\rs$
is precomputed data, this is solved by an outer loop (\code{bayes.estimate}, \S\ref{sec:Tloop}).
\code{softmin_T_bps} overrides the loop with one fixed $T$ for all segments.
\item \textbf{Numerics} (\code{model.softmin_rate}): with $z_j=-q_j/T$ and $m=\max_j z_j$, the code
computes $\rs=-T\big(m+\log\sum_je^{z_j-m}\big)$. The largest term becomes $e^0=1$, so the sum cannot
underflow to zero. Without the shift, $e^{-600/1}$ is $0$ in double precision. Absent lenders
(NaN) are set to $z=-\infty$ and contribute $0$. A row with no lenders is an error.
\item \textbf{Period convention} (\code{panel._per_period}): all rates are day-weighted averages over the
period $(t-1,t]$. A change log is expanded to a daily step function, forward-filled from each
effective date, and averaged over the days of each period. The first period is one median spacing
long, and days before a lender's first change are unknown and skipped. Why the convention matters:
suppose a move of $\Delta$ takes effect a fraction $\theta$ of the way through a period. Its day-weighted
average moves by $(1-\theta)\Delta$ in that period, while an end-of-period rate moves by $\Delta$. Mixing
the two conventions creates a spurious one-period gap of $\theta\Delta$ around every common move.
\item \emph{Approximation:} $\rs$ is the soft-min of period-average rates. It is not the period
average of daily soft-mins. Because $\rs$ is concave (its Hessian is $-(\mathrm{diag}\,w-ww^\top)/T\preceq0$),
the two differ only when lenders reprice within a period, and the difference is small for weekly data.
\end{itemize}

\takeaway{Borrowers compare us with the whole field, but mostly with the cheapest lenders. The
benchmark $\rs$ is a ``soft cheapest rate'': it follows the keenest lenders, ignores the uncompetitive ones,
and falls a little whenever another option appears. How soft it is depends on how price-sensitive
borrowers are, so it is solved together with the sensitivity.}

%=====================================================================
\section{The price gap $x$}\label{sec:x}

\subsection{The unlagged gap}
\code{panel.add_gap} sets, per segment and week,
\begin{equation}\label{eq:gap}
v_{kt}=H_{kt}+\kappa\,D_{kt}-\big(\rs_{kt}-\rs_{k1}\big).
\end{equation}
Why it is \eqref{eq:binary}: our perceived rate is $r_{us,kt}=r_{us,k1}+H_{kt}+\kappa D_{kt}$, so
$r_{us}-\rs=v_{kt}+(r_{us,k1}-\rs_{k1})$. The bracket is a segment constant and is absorbed by
$\alpha_k$ (\S\ref{sec:alpha}). Without a competitor file the last term is zero, competitor moves are
invisible, and the run warns.

\paragraph{What $\kappa$ means.} A 1\,bp change in our average discretionary discount moves the
perceived price by $\kappa$\,bp. Discounts reach only some borrowers, are negotiated after the
comparison stage, and are less visible than the carded rate. So $\kappa<1$ is the average
``visibility'' of the discount. It is fixed at $0.5$ and not estimated: with about ten moves, few of
them discretionary, the data cannot separate $\kappa$ from $\beta$ (only the product $\beta\kappa$
enters for a discretionary move). The sensitivity runs refit at $\kappa=0.3$ and $0.7$
(\S\ref{sec:check}).

\subsection{Partial adjustment}
Customers and brokers notice price changes gradually. \code{model.lag} applies, per segment,
\[
p_t=p_{t-1}+\phi\,(v_t-p_{t-1}),\qquad p_0:=v_1,\qquad \phi=1-2^{-1/h},
\]
and $x_{kt}=p_{kt}$. There is one half-life $h$, shared by all segments and by all three components
of $v$. The earlier design had separate half-lives for headline and discretionary moves, chosen with
$\kappa$ by a grid search; about ten moves cannot support three timing parameters.
\begin{lemma}\label{lem:lag}
(i) $x_t=v_1+\sum_{j=1}^{t}\phi(1-\phi)^{t-j}(v_j-v_1)$. In matrix form, with the lower-triangular
kernel $W_{tj}(h)=\phi(1-\phi)^{t-j}\one[j\le t]$,
\begin{equation}\label{eq:kernel}
x_k=v_{k1}\one+W(h)\,(v_k-v_{k1}\one).
\end{equation}
(ii) After a step of size $\Delta$ in $v$ at week $\tau_0$, the part felt $\tau$ weeks in (counting
the week of the step as $\tau=1$) is $\Delta\,\big(1-(1-\phi)^{\tau}\big)$, so half is felt after $h$ weeks.
(iii) The average part felt over the first $N$ weeks is
\[
\bar f_N(h)=1-\frac{(1-\phi)\big(1-(1-\phi)^N\big)}{N\phi}.
\]
\end{lemma}
\begin{proof}
(i) Put $d_t=v_t-v_1$ and $p'_t=p_t-v_1$. Then $p'_t=(1-\phi)p'_{t-1}+\phi d_t$ with $p'_0=0$, and
iterating gives $p'_t=\sum_{j\le t}\phi(1-\phi)^{t-j}d_j$. (ii) The remaining gap obeys
$g_\tau=(1-\phi)g_{\tau-1}$ with $g_0=\Delta$; $(1-\phi)^h=\tfrac12$ defines $\phi$. (iii) Average
(ii) over $\tau=1,\dots,N$ and sum the geometric series.
\end{proof}
At $h=2$: $\phi=0.293$ and $\bar f_{12}=1-0.707\times0.984/(12\times0.293)=0.80$. \emph{Only about
80\% of a move is felt on average over its first 12 weeks.} This is why \S\ref{sec:forecast}
forecasts a path and does not apply the steady-state effect. Starting at $p_0=v_1$ assumes the
panel begins in steady state.

\paragraph{Why the kernel.} $h$ is a model parameter, with prior
$h\sim\mathrm{LogNormal}(\log2,\,0.5)$ truncated to $[0.5,12]$ weeks, so the lag must be a
differentiable function of $h$ inside the PyMC graph. A recursion over $t$ would need a scan. The
kernel \eqref{eq:kernel} is one matrix product: $v$ is precomputed data, and
\code{bayes.lag_kernel_pt} builds $W$ as $\phi\exp\{(t-j)\log(1-\phi)\}$ on the lower triangle. A unit
test checks \eqref{eq:kernel} against the recursion for $h\in\{0.5,1,2,5,12\}$.

\subsection{Unseen competitor discounts: measurement error}
We see competitors' carded rates and \code{adj}, but not their discretionary discounts. So the
observed gap is $x^{\rm obs}=x+m$, where $m$ is the (lagged, weighted) error in $\rs$; to first order
$m_{kt}\approx\sum_j w_{jkt}\,\Delta\text{disc}_{jkt}$. In a linear regression $y=a-\beta x+e$ with
classical error $m\perp(x,e)$,
\[
\plim\hat\beta=\beta\,\frac{\Var(x)}{\Var(x)+\Var(m)}\in(0,\beta],
\]
and the same attenuation operates in the present model, where $y$ is the logit share net of the
natural level. The persistent part of $m$ is absorbed by $\alpha_k$ and its smooth part by $g_t$
and $\nu_{kt}$; what attenuates $\beta$ is the part of $m$ that moves like a price move. This is
large in practice: discounts on advertised rates have moved by tens of bp within a year, more than
most single repricings. The error is not classical if lenders trade headline against discount; then
the bias can go either way. The only defence inside the model is \code{adj_bps}.

\takeaway{$x$ is how much dearer we look than when the data starts, relative to the market, after
allowing customers time to notice. One number, the half-life, sets the speed, and the data update
it. At a two-week half-life about 80\% of a move is felt on average over its first 12 weeks.
Competitor discounts we cannot see blur $x$ and make customers look less price-sensitive than they are.}

%=====================================================================
\section{Natural share}\label{sec:alpha}

\subsection{The level}
Equation \eqref{eq:binary} is $\logit\pi_{kt}=a_{us,kt}-\beta_k(r_{us}-\rs)$. Everything that is not
price sits in $a_{us,kt}$. The model writes it as a level that drifts:
\begin{equation}\label{eq:level}
\logit\pi_{kt}=\underbrace{\alpha_k+g_t+\nu_{kt}+J_{kt}}_{\ell_{kt}}\;-\;\beta_k\,x_{kt}(h),
\end{equation}
and $\Lambda(\ell_{kt})$ is the \emph{natural share}: the share we would have at the price gap of
week 1. The pieces (\code{bayes.build_model}):
\begin{align*}
\alpha_k&\sim N\big(\logit s_{0k},\,1\big), &&\text{segment constant;}\\
g_t&=g_{t-1}+\sigma_g\,\xi_t,\quad g_1=0,\ \xi_t\sim N(0,1), &&\sigma_g\sim\mathrm{HalfNormal}(0.05);\\
\nu_{kt}&=\nu_{k,t-1}+\sigma_v\,\zeta_{kt},\quad \nu_{k1}=0,\ \textstyle\sum_k\zeta_{kt}=0, &&\sigma_v\sim\mathrm{HalfNormal}(0.05);\\
J_{kt}&=\textstyle\sum_{e}j_e\,\one[k\in S_e,\ t\ge t_e], &&j_e\sim N(0,\,0.30^2).
\end{align*}
Turnaround time, credit policy, service and brand are in $\ell$ by design. There is no turnaround
covariate. An optional matrix of standardised covariates $z_{kt}$ with coefficients
$\gamma\sim N(0,0.2^2)$ can be added through \code{Config.covariates}; the default is none.

\paragraph{Zero-sum segment innovations.} $\zeta_{\cdot t}\sim\mathrm{ZeroSumNormal}$: a standard
normal vector projected onto $\{\zeta:\one^\top\zeta=0\}$, with covariance $I-\one\one^\top/K$.
Without the constraint, a common shock could be put in $g_t$ or spread over the $\nu_{kt}$, and the
posterior would have a flat direction. With it, $g_t$ is exactly the cross-segment mean of the
drift. The variance of one segment's weekly drift is $\sigma_g^2+\sigma_v^2(1-1/K)$.

\paragraph{Non-centred form.} The model samples the standardised innovations $\xi,\zeta$ and
builds $g=\sigma_g\,\mathrm{cumsum}(\xi)$, $\nu=\sigma_v\,\mathrm{cumsum}(\zeta)$. Weekly noise
\eqref{eq:sampling} is several times the weekly drift, so each state is only weakly informed by its
own week. In that regime the centred form has a funnel between $\sigma$ and the states and the
non-centred form does not.

\subsection{Why a random walk and not week dummies}\label{sec:rw}
The earlier design used free week effects $\delta_t$. Let a bank-wide move put the same $x_t$ in
every segment.
\begin{proposition}\label{prop:common}
With free week effects and a common slope $\beta$, the likelihood depends on $(\delta_t,\beta)$ only
through $\delta_t-\beta x_t$. A bank-wide move carries no information about $\beta$. With
segment slopes $\beta_k=\bar\beta+b_k$, it identifies the differences $b_k$ but not the level
$\bar\beta$.
\end{proposition}
\begin{proof}
$\logit\pi_{kt}=\alpha_k+\delta_t-\beta_kx_t=\alpha_k+(\delta_t-\bar\beta x_t)-b_kx_t$. The map
$(\delta_t,\bar\beta)\mapsto(\delta_t+cx_t,\ \bar\beta+c)$ leaves every $\pi_{kt}$ unchanged.
\end{proof}
About half of our moves are bank-wide, so free dummies discard about half of a small sample. A
random walk does not: its prior says $g_t-g_{t-1}$ is small, so a week effect cannot follow a
price step for free.
\begin{proposition}[Information in a bank-wide move]\label{prop:info}
Let the common component obey $y_t=g_t-\beta x_t$ with $g$ a random walk of innovation sd
$\sigma_g$, observed without sampling noise, and let $x$ be a step $\Delta$ at one week, lagged with
speed $\phi$. Then the Fisher information for $\beta$ is
\[
\mathcal I(\beta)=\frac{\Delta^2}{\sigma_g^2}\cdot\frac{\phi}{2-\phi},
\qquad\text{so}\qquad \mathrm{sd}(\hat\beta)\ \ge\ \frac{\sigma_g}{|\Delta|}\sqrt{\frac{2-\phi}{\phi}} .
\]
\end{proposition}
\begin{proof}
First differences are independent: $y_t-y_{t-1}=\sigma_g\xi_t-\beta\,(x_t-x_{t-1})$, and by
Lemma~\ref{lem:lag}(ii) $x_t-x_{t-1}=\Delta\phi(1-\phi)^{\tau-1}$ in week $\tau$ of the move. So
$\mathcal I=\sigma_g^{-2}\sum_{\tau\ge1}\Delta^2\phi^2(1-\phi)^{2(\tau-1)}=\sigma_g^{-2}\Delta^2\phi^2/(1-(1-\phi)^2)$.
\end{proof}
At $h=2$ the factor $\sqrt{(2-\phi)/\phi}$ is $2.4$. With $\sigma_g=0.015$ (about the value on the
simulated sample) one 10\,bp bank-wide move gives $\mathrm{sd}(\hat\beta)\ge0.0036$ per bp, about
$0.3\%$/bp in uplift terms; $m$ such moves divide it by $\sqrt m$. Sampling noise adds to this, more in thin segments.
Two readings of the formula matter. First, information grows with $\phi$: a slow response
($\phi\to0$) is indistinguishable from slow drift, so $h$ and $\beta$ are only jointly identified
when $h$ is long. Second, everything rests on $g$ being smooth:

\paragraph{The assumption.} \emph{No other abrupt bank-wide shock lands in the same weeks as a
bank-wide price move, unless it is logged.} An unlogged shock of size $\delta$ in the week of a
move $\Delta$ biases $\beta$ by about $-\delta/\Delta$. \code{Config.time_effect = "fe"} replaces
$g$ by sum-to-zero week effects with a $N(0,0.5^2)$ prior; it exists for the sensitivity run that
shows how much of the estimate rests on bank-wide moves (\S\ref{sec:check}).

\subsection{Jumps for logged non-price changes}
Each event $e$ of type \code{other} gets one permanent shift $j_e$, shared by its segments $S_e$,
from its week $t_e$ on. A price move in the same segments and week then competes with $j_e$: the
long-run change in logit share is $j_e-\beta_k\Delta$, and the data separate the two only through
the \emph{shape} of the response (an immediate jump against a phased-in move), which is weak. So
such a price move teaches the model little. That is the intent: its share response cannot be put
down to price. Price events add nothing to \eqref{eq:level}; their effect is carried by $x$.

\subsection{Moves made because share had slipped}
Suppose we cut price after weak volume, as \code{simulate --reactive} does: $x_t$ depends on past
shares. A regression with a fixed level treats the slipped level as noise, and since the cut is
correlated with it, $\hat\beta$ is biased towards zero (\S\ref{sec:diag-est} derives the size). In
\eqref{eq:level} the level is a parameter. Write the joint density of the data as
\[
p(L,x\mid\theta)=\prod_t p\big(L_{\cdot t}\mid \ell_{\cdot t},x_{\cdot t},\beta,h,c\big)\;
p\big(x_{\cdot t}\mid L_{\cdot,<t},x_{\cdot,<t}\big).
\]
If the pricing rule depends only on what has been observed, its factor does not involve $\theta$ and
drops out of the posterior: conditioning on $x$ is valid, and a cut is measured against where the
level had drifted to. Two limits remain. A move made on information that is not yet in share
(a forward view of a competitor's plan, say) is still confounded. And the drifting level costs
precision: the share movement after a cut can be partly attributed to drift, which is what
Proposition~\ref{prop:info} quantifies.

\takeaway{Share moves for many reasons besides price. The model keeps a running estimate of the
share we would have had at an unchanged price gap, made of a smooth bank-wide trend, a smooth trend
per segment and a one-off jump at each logged non-price change. A price cut is then judged against
where share had drifted to. Bank-wide price moves can be used only because the bank-wide trend is
assumed smooth, so log any other bank-wide shock.}

%=====================================================================
\section{Sensitivity: prior and pooling}\label{sec:prior}

\subsection{Uplift per bp}
The business thinks in uplift: the fractional rise in volume per 1\,bp cut. From
$s=\Lambda(\ell-\beta x)$, $\partial\log s/\partial x=-\beta(1-s)$, so at the baseline share
\[
u_k=\beta_k\,(1-s_{0k}),\qquad \beta_k=\frac{u_k}{1-s_{0k}}
\]
(\code{model.beta_from_uplift}). $s_{0k}=\sum_tL_{kt}/\sum_tM_{kt}$ is a fixed constant used only
for this conversion; it makes $u_k$ comparable across segments of different share.

\subsection{The hierarchy}
\begin{align*}
\log u_k&=\mu_u+a_{\rm inv}\big(\one[\text{INV}]-\tfrac12\big)+a_{\rm io}\big(\one[\text{IO}]-\tfrac12\big)+c^{\rm lvr}_{b(k)},\\
\mu_u&\sim N(\log0.03,\ 0.70^2),\qquad a_{\rm inv},a_{\rm io}\sim N(0,0.30^2),\\
c^{\rm lvr}&\sim\mathrm{ZeroSumNormal}(\sigma_{\rm lvr})\ \text{over the 5 bands},\qquad \sigma_{\rm lvr}\sim\mathrm{HalfNormal}(0.30).
\end{align*}
The $\pm\tfrac12$ coding makes $\mu_u$ the log uplift of an average segment, and the zero-sum
constraint does the same across bands. There is no segment-level residual by default: twenty
$u_k$ are functions of seven free numbers ($\mu_u$, two contrasts, four band effects). That is the
pooling. With about ten moves, most segments have no move of their own, and a segment's $u_k$ is
learned from every segment that shares its purpose, repayment type or band.
\code{Config.segment_residual} adds $\varepsilon_k\sim N(0,\sigma_\varepsilon^2)$,
$\sigma_\varepsilon\sim\mathrm{HalfNormal}(0.15)$. Without it, a segment that truly departs from
the additive pattern has an interval that is too narrow.

\subsection{What the prior says}
For an average segment $u=e^{\mu_u}$ is lognormal: median $3\%$/bp, and
\[
\text{80\% interval}=0.03\,e^{\pm1.2816\times0.70}=[1.22\%,\ 7.36\%],\qquad
\Pr(u>8\%)=1-\Phi\!\Big(\frac{\log(8/3)}{0.70}\Big)=0.08 .
\]
So $8\%$/bp, the earlier assumption, sits at about the 92nd percentile: possible, not assumed. For a
particular segment the offsets add variance,
$\Var(\log u_k)=0.70^2+2\times\tfrac14\times0.30^2+\E[\sigma_{\rm lvr}^2]\,(1-\tfrac15)=0.49+0.045+0.072$,
a standard deviation of $0.78$; \code{estimate} draws from the prior and prints the implied
interval for a mid-range segment (about $1.1$ to $8.3\%$/bp).

\paragraph{Where the numbers come from.} Published lender-level estimates, none Australian and all
on transacted prices, sit at roughly $1$ to $2.5\%$ of share per bp. The median is set above that
range because the lender is small and price-focused. That shift is a judgment and is not sourced.
The prior median and spread are \code{prior_u_median} and \code{prior_u_log_sd}; the sensitivity
run refits at medians of 2, 4 and 8\%/bp.

\subsection{How much of the answer is prior}\label{sec:priorshare}
\code{report.prior_share} is
\[
\texttt{prior\_share}_k=\min\Big\{1,\ \frac{\Var_{\rm post}(\log u_k)}{\Var_{\rm prior}(\log u_k)}\Big\},
\]
with the prior variance taken from prior draws of the same model. In the normal--normal case, with
prior variance $\sigma_0^2$ and data variance $\sigma_d^2$, the posterior variance is
$(\sigma_0^{-2}+\sigma_d^{-2})^{-1}$ and the ratio is
$\sigma_0^{-2}/(\sigma_0^{-2}+\sigma_d^{-2})$: the share of posterior precision that came from the
prior, which is also the weight of the prior mean in the posterior mean. Near 1, the result is mostly
prior. Because of pooling, a segment with no move of its own can still have a low value.

\takeaway{Sensitivity starts from a documented range (middle 3\% per bp, 80\% range 1.2 to 7.4),
not a single figure, and the old 8\% is inside it but unlikely. Segments that share a purpose,
repayment type or LVR band inform each other. Each segment's output says how much of its answer is
still the starting range.}

%=====================================================================
\section{Likelihood and estimation}\label{sec:fit}

\subsection{Beta-binomial counts}
\begin{equation}\label{eq:lik}
L_{kt}\mid M_{kt}\sim\mathrm{BetaBinomial}\big(M_{kt},\ \pi_{kt}c,\ (1-\pi_{kt})c\big),\qquad
c\sim\mathrm{LogNormal}(\log1000,\ 1.5).
\end{equation}
Equivalently $\tilde\pi\sim\mathrm{Beta}(\pi c,(1-\pi)c)$ and $L\mid\tilde\pi\sim\mathrm{Bin}(M,\tilde\pi)$. Then
\[
\E L=M\pi,\qquad \Var L=M\pi(1-\pi)\Big(1+\frac{M-1}{c+1}\Big).
\]
The bracket is the overdispersion: week-to-week wobble in share beyond binomial sampling and beyond
what the random walks carry. For $c\to\infty$ it is plain binomial. At $c=1000$ a segment with
$M=4000$ has five times the binomial variance, and one with $M=60$ has almost none: extra-binomial
noise matters where volumes are large, and sampling noise dominates where they are thin. Compared
with the earlier regression on $\logit s$ this handles $L=0$, weights every week correctly, and
puts the noise where it arises.

\subsection{The posterior}
The unknowns are $\theta=(\mu_u,a_{\rm inv},a_{\rm io},c^{\rm lvr},\sigma_{\rm lvr},\alpha,\sigma_g,\xi,\sigma_v,\zeta,j,h,c)$:
for 20 segments and 104 weeks about 2{,}200 numbers, nearly all of them weekly innovations. The
posterior is \eqref{eq:lik} over the observed cells, times the priors of
\S\S\ref{sec:x}--\ref{sec:prior}, with $\pi$ from \eqref{eq:level} and $x$ from \eqref{eq:kernel}.
It has no closed form. \code{bayes.fit} samples it with NUTS (nutpie if installed, else numpyro, else PyMC's
sampler): 4 chains, 1{,}000 warm-up and 2{,}000 kept draws each, \code{target_accept} $=0.95$, a fixed seed. Every
reported quantity is a function of the draws, and every interval is the 10th to 90th percentile
across them.

\paragraph{Diagnostics.} After each full fit the run reports the largest $\hat R$ over all
parameters, the smallest bulk effective sample size over the key parameters
($\mu_u$, $a_{\rm inv}$, $a_{\rm io}$, $h$, $c$, $u_k$) and the number of divergences. It writes
\texttt{UNRELIABLE RUN} if $\hat R>1.01$ anywhere, an ESS is below 400, or there is any divergence.
With 1{,}000 draws per chain the largest of some 2{,}200 $\hat R$ values sat at the $1.01$ limit and
tripped it in several of a handful of trial runs, which is why the default is 2{,}000.

\subsection{Solving the temperature}\label{sec:Tloop}
The choice model requires $T_k=1/\beta_k$ (\S\ref{sec:choice}), but $\rs$, and through it $x$, is
precomputed with a given $T$. The consistent pair is a fixed point of
$T\mapsto1/\hat\beta(T)$, where $\hat\beta(T)$ is the posterior median fitted on the benchmark built
with $T$. \code{bayes.estimate}:
\begin{enumerate}
\item starts at $T_k=(1-s_{0k})/u_{\rm prior}$, i.e.\ $1/\beta_k$ at the prior median;
\item runs a short fit (2 chains, 500 warm-up, 500 draws);
\item updates with damping, $T^{\rm new}_k=\sqrt{T^{\rm old}_k/\hat\beta_k}$, clipped to $[2,60]$\,bp;
\item rebuilds $\rs$ and $v$, and repeats until $\max_k|T^{\rm new}_k/T^{\rm old}_k-1|<5\%$ or 6 iterations;
\item runs the full fit at the final $T$.
\end{enumerate}
In logs, with $\tau=\log T$ and $F(\tau)=-\log\hat\beta(e^{\tau})$, step 3 is
$\tau^{\rm new}=\tfrac12\big(\tau+F(\tau)\big)$. Near the fixed point the error is multiplied by
$\tfrac12(1+F')$ each round, so the damped iteration converges whenever $-3<F'<1$, where the undamped
one needs $|F'|<1$. In practice $F'$ is small and negative or positive ($\hat\beta$ moves little
with $T$), so the error roughly halves each round. Note what the stopping rule measures: a change
below 5\% means the remaining error is about that size too. The path of $T$ per iteration is
written to \code{T_path.csv}, and the run warns loudly if the loop did not converge.

\paragraph{Weak joint identification.} $T$ changes how much weight $\rs$ gives to lenders that are
not the cheapest. The data distinguish values of $T$ only through competitor moves away from the
cheapest rate, so $T$ and $\beta$ are jointly weakly identified and the loop mostly imposes the
theory's restriction. The uncertainty in $T$ is not in the reported intervals. The sensitivity run
refits at $T/2$ and $2T$.

\takeaway{One model is fitted in one go. The fit produces thousands of equally plausible versions
of the whole model, and every reported range is read off them. The softness of the benchmark is
solved in a loop until it agrees with the estimated sensitivity. The run checks its own numerical
health and says so plainly when a fit should not be used.}

%=====================================================================
\section{Forecast and the Excel export}\label{sec:forecast}

\subsection{The multiplier}
Hold the market total $M$ and conversion fixed. With $s(x)=\Lambda(a-\beta x)$ and
$s_0=s(0)=\Lambda(a)$, we have $e^{-a}=(1-s_0)/s_0$, so
\[
\frac{A(x)}{A_0}=\frac{s(x)}{s_0}=\frac{1+e^{-a}}{1+e^{-a}e^{\beta x}}
=\boxed{\frac{1}{s_0+(1-s_0)e^{\beta x}}}
\]
(\code{model.multiplier}).
Properties: $m(0)=1$, $m'(0)=-\beta(1-s_0)=-u$, and
$m''(0)=\beta^2(1-s_0)(1-2s_0)$. So for $s_0<\tfrac12$ the curve is convex: a cut helps by more than
an equal rise hurts. As $x\to-\infty$, $m\to1/s_0$ (the whole market). What an uplift means for a
move is therefore not $u\Delta$: at $s_0=10\%$, a 10\,bp cut multiplies volume by $2.13$ at $8\%$/bp,
$1.34$ at $3\%$/bp and $1.22$ at $2\%$/bp, in the long run.

\subsection{The path}
\code{forecast.forecast} takes the segments, our headline and discretionary change, optional
competitor changes and a horizon, and works draw by draw. The scenario changes the unlagged gap
from week 1 by
\[
\Delta x=\Delta H+\kappa\,\Delta D-\big(\rs(q')-\rs(q)\big),
\]
with $\rs(q')$ recomputed at the latest effective rates plus the competitor changes, in
\code{forecast.benchmark_shift}. For small moves $\rs(q')-\rs(q)\approx\sum_jw_j\Delta q_j$, by
Proposition~\ref{prop:softmin}(f). Let $x_n$ be the lagged gap in the last week and $v_n$ the
unlagged one. By Lemma~\ref{lem:lag}(ii), holding the gap at its target,
\begin{align*}
\text{baseline:}\quad x^{0}_{n+\tau}&=x_n+(v_n-x_n)\big(1-(1-\phi)^\tau\big),\\
\text{scenario:}\quad x^{1}_{n+\tau}&=x_n+(v_n+\Delta x-x_n)\big(1-(1-\phi)^\tau\big).
\end{align*}
The baseline is not flat when a recent move has not finished phasing in ($x_n\ne v_n$), which is
why it is computed and not assumed. Shares are $s^{i}_{\tau}=\Lambda(\ell_{n}-\beta x^{i}_{n+\tau})$
with $\ell_n$ the last natural-share level, and
\[
m_\tau=\frac{s^1_\tau}{s^0_\tau},\qquad
m_{1:N}=\frac{\sum_{\tau=1}^{N}s^1_\tau}{\sum_{\tau=1}^{N}s^0_\tau}\quad(N=12).
\]
From a settled baseline ($x_n=v_n$) $s^0$ is constant, $m_\tau=m\big(\Delta x\,(1-(1-\phi)^\tau)\big)$ and
$m_{1:N}$ is the plain average of the weekly multipliers. \code{share_now} in \code{segments.csv} is
$\Lambda(\ell_n-\beta x_n)$. The steady-state multiplier $m(\Delta x)$ is reported separately as the
long-run effect. By Lemma~\ref{lem:lag}(iii) it overstates the first 12 weeks by about a fifth of the effect.

\subsection{Two intervals}
\begin{itemize}
\item \textbf{Price-effect interval.} $\ell$ is held at $\ell_n$. The spread across draws reflects
$\beta$, $h$ and the estimate of $\ell_n$ only. This is the business condition ``no competitor
response and no market movement''. Lodgements are $\bar M_k s_\tau$ with $\bar M_k$ the trailing
13-week mean market volume, and applications are lodgements times the conversion ratio.
\item \textbf{Full predictive interval.} The random walks are simulated forward,
$\ell_{k,n+\tau}=\ell_{kn}+\sigma_g\sum_{i\le\tau}\xi_i+\sigma_v\sum_{i\le\tau}\zeta_{ki}$ with zero-sum $\zeta$,
and counts are drawn from \eqref{eq:lik} at $\bar M_k$. Its multiplier is the simulated scenario
total divided by the posterior-median baseline total. With \code{time_effect = "fe"} there is no
bank-wide trend to project and that part is held flat.
\end{itemize}
The drift term is worth sizing. The average of a random walk over its first $N$ steps has variance
$\sigma^2(N+1)(2N+1)/(6N)$, which is $4.5\sigma^2$ at $N=12$: a standard deviation of about
$2.1\times$ the weekly drift, relative to volume (times $1-s$).

\paragraph{Groups.} For a group of segments, lodgements are summed across segments \emph{within each
draw} before percentiles are taken, so common uncertainty ($\mu_u$, $h$, $g$) does not average
away and segment-specific noise does.

\subsection{The workbook}
\code{excel.export_excel} writes live formulas (openpyxl). \code{segments} holds $s$ = \code{share_now},
$\beta$ at p10/p50/p90, $T_k$, the median $h$, conversion and \code{prior_share}. \code{benchmark}
holds each lender's latest effective rate, an editable scenario rate, and, now and under the
scenario,
\[
w_j=\frac{e^{-(q_j-q_{\min})/T}}{\sum_le^{-(q_l-q_{\min})/T}},\qquad
\rs=q_{\min}-T\log\sum_je^{-(q_j-q_{\min})/T}.
\] \code{scenario} takes $\Delta H$ and $\Delta D$ per segment, computes
$\Delta x$, and for $\tau=1,\dots,12$
\[
m_\tau=\frac{1}{s+(1-s)\exp\!\big\{\beta\,\Delta x\,\big(1-(1-\phi)^{\tau}\big)\big\}},\qquad \phi=1-2^{-1/h},
\]
then averages over the 12 weeks at $\beta$ p10, p50 and p90, and shows the long-run multiplier. The
p10--p90 range is labelled ``80\% range from sensitivity only'': it starts from a settled baseline,
fixes $h$ at its median and $s$ at its median, so it is narrower than \code{forecast}. A test
evaluates the workbook's formulas and checks them against \code{forecast} (price-effect version,
settled baseline, $h$ at its median) to within $0.5\%$.

\takeaway{A forecast plays every plausible version of the model forward 12 weeks with and without
the change. A move takes weeks to sink in, so the 12-week effect is smaller than the long-run one.
Two ranges are given: one for the price effect alone, and a wider one that also allows for drift
and chance. The Excel sheet reproduces the first with formulas you can see.}

%=====================================================================
\section{Accuracy, diagnostics and validation}\label{sec:check}

\subsection{What accuracy the data can deliver}
The business target is 12-week applications within $\pm5\%$, and share within $\pm0.3$ points, at
80\% confidence. An 80\% interval of $\pm5\%$ needs a total relative standard deviation of
$5\%/1.2816=3.9\%$. Three sources add in quadrature:
\begin{center}
\footnotesize
\begin{tabular}{@{}lll@{}}
\toprule
Source & Size (one sd, relative) & Controlled by\\
\midrule
Sensitivity & $|\Delta x|\ \bar f_{12}\ \mathrm{sd}(u)$ & number and cleanliness of past moves\\
Sampling & $\sqrt{(1-s)\,/\,\textstyle\sum_{\tau\le12}L_\tau}$, times overdispersion & segment volume\\
Drift in natural share & $\approx2.1\,(1-s)\sqrt{\sigma_g^2+\sigma_v^2}$ & nothing in the model\\
\bottomrule
\end{tabular}
\end{center}
Consequences. (i) To hold $\pm5\%$ from sensitivity alone, $\mathrm{sd}(u)$ must be under
$0.039/(|\Delta x|\bar f_{12})$: about $1.0$, $0.5$ and $0.2\%$/bp for 5, 10 and 25\,bp moves. (ii)
Sampling alone fits inside the target only if $\sum L\ge(1-s)/0.039^2\approx590$ over 12 weeks at
$s=10\%$, about 50 own lodgements a week; below that (\code{thin_volume}) the target is meaningful
only for groups. (iii) $\pm0.3$ points is tighter than $\pm5\%$ of applications when share is
above $6\%$ and looser below. (iv) Drift is excluded by the ``no market movement'' condition, which
is a counterfactual: it is never observed, and can be tested only by backtests that feed in the
competitor rates and market volumes that occurred.

\code{forecast.accuracy_table} computes this from the fitted model and not from the approximation:
for each segment and group (all, by purpose, by repayment type, by LVR band) and each own headline
move of $\pm5$, $\pm10$, $\pm25$\,bp, the half-width of the 80\% interval of 12-week applications
relative to its median, and of 12-week share in points, for both interval versions, with
\code{meets_5pct} and \code{meets_0p3pp}. Nothing is tuned to make it pass.

\subsection{The raw-move check}
\code{diagnostics.raw_move_check} uses no model and is meant to be read before any fit. For each
logged own price move and affected segment it reports the change in the unlagged gap
(mean over weeks $3$ to $10$ after, less mean over the 4 weeks before), mean weekly applications and
share in the same two windows, their ratio, and the ratio implied by an uplift $u$:
$m(\Delta v)$ with $\beta=u/(1-s_{\rm before})$, for $u=8\%$ and $3\%$/bp. Under $8\%$/bp a 10\,bp
cut should roughly double applications; if the raw data show nothing like that, no model is needed
to doubt the figure. The comparison is crude: it ignores drift and takes the benchmark at the
prior-median temperature.

\subsection{The old estimators, as diagnostics}\label{sec:diag-est}
\code{diagnostics.estimators} reports two estimates from the earlier design beside the model's
$u_k$. Neither feeds the estimate. Both use $y=\logit s$, so segments with a zero-share week are left out.

\paragraph{Two-way fixed effects} (\code{fe.fit_twfe}): OLS of
$y_{kt}=\alpha_k+\delta_t-\beta_kx_{kt}+e_{kt}$ with HC1 errors, $x$ at the posterior-median $h$.
Its known biases:
\begin{itemize}
\item By Proposition~\ref{prop:common}, $\delta_t$ absorbs every bank-wide move.
\item \emph{Reactive repricing.} Let the true level be a random walk and the fixed effect capture
only its mean, so $e_{kt}=(\ell_{kt}-\bar\ell_{k\cdot})+\varepsilon_{kt}$, and let
$x_{kt}=x^{\rm exo}_{kt}+\psi\,\ell_{k,t-1}$ with $\psi>0$ (a low level leads to a lower price). For
a common slope, after partialling out the fixed effects (tildes),
\[
\plim\hat\beta=\beta-\frac{\Cov(\tilde x,\tilde e)}{\Var(\tilde x)}
=\beta-\frac{\psi\Var(\tilde\ell)}{\Var(\tilde x^{\rm exo})+\psi^2\Var(\tilde\ell)}<\beta ,
\]
because for a random walk $\Cov(\ell_{t-1},\ell_t)=\Var(\ell_{t-1})$. The bias is towards zero, it
needs persistence (for an AR(1) level with coefficient $\varrho$ the covariance is
$\psi\varrho\Var(\ell)$), and since $\Var(\tilde\ell)$ grows with $n$ it does not shrink with more data.
\item \emph{Standard errors.} HC1 assumes the products $\tilde x_t\tilde e_t$ are uncorrelated over
$t$. With AR(1) coefficients $\varrho_x,\varrho_e$ the true variance is larger by
$(1+\varrho_x\varrho_e)/(1-\varrho_x\varrho_e)$: a factor $6.1$ at $0.9$ and $0.8$, so the true SE
is about $2.5\times$ the reported one.
\item Unseen competitor discounts attenuate it (\S\ref{sec:x}).
\end{itemize}

\paragraph{Event studies} (\code{events.event_study}): for a move in segment $k$,
$\hat\beta_e=-dy/dx$ with $dy$ and $dx$ the change from the 4 weeks before to the 8 weeks after, less
the same change in segments with no move in the window. Its known biases: a bank-wide move has no
control segments and is dropped; a move made after share had slipped is biased towards zero for the
reason above; a move with a logged non-price change in the window is excluded; and with a handful
of usable events it is noisy. The table reports the precision-weighted mean per segment and the
number of events behind it.

\subsection{Posterior predictive checks}
\code{diagnostics.ppc} simulates lodgement histories from the fitted model (400 draws) and compares,
per segment, two statistics of weekly share with their simulated 80\% range: the variance, which
tests the noise model \eqref{eq:lik}, and the lag-1 autocorrelation, which tests the dynamics. An
observed value far outside its range says that the model wobbles or persists differently from the data.

\subsection{Backtests}
All refits hold $T$ at the base fit's converged value (\code{validate.py}). A cell left out of the
likelihood still has its latent level, which the random-walk prior carries through the gap, so a
masked fit predicts masked cells exactly as a forecast would, with the uncertainty of the drift.
\begin{itemize}
\item \textbf{Leave one move out} (\code{validate events}). For each own price move, the 12 weeks
from the move in the affected segments are masked. The refit predicts those weeks' lodgements from
the actual market volumes and the actual gap path. Reported: predicted 12-week total
(p10, p50, p90) against actual, the relative error, and the coverage of the 80\% interval. This is
the test of the counterfactual target in the only form available.
\item \textbf{Quiet periods} (\code{validate quiet}). Rolling-origin 12-week forecasts from origins
every 13 weeks after week 52, with everything after the origin masked, scored on segments with no
own move in the window. The error here is the floor from drift and sampling.
\item \textbf{Placebo} (\code{validate placebo}). A fake 10\,bp step at a random quiet date and
segment enters with its own coefficient (prior $N(0,0.1^2)$ per bp), lagged with the same kernel.
Its 80\% interval should contain zero about 80\% of the time over the 10 repetitions.
\item \textbf{Sensitivity} (\code{sensitivity}). Refits at prior medians 2, 4 and 8\%/bp; $T$ halved
and doubled; $\kappa=0.3$ and $0.7$; week dummies in place of the smooth trend; jumps off. One table
gives $u$ (p10, p50, p90), as the lodgement-weighted mean of $u_k$ within each draw, overall and by
group, beside the base run.
\end{itemize}

\subsection{Acceptance on simulated data}\label{sec:accept}
\code{simulate.py} generates data from the model with known $u_k$: 20 segments, several thin,
about 12 own moves split between bank-wide and segment-specific, RBA-like steps passed on by every
lender after its own delay, competitor specials, logged non-price jumps, and optionally reactive
repricing. The truth is stated as the model defines it, $u_k=\beta_k(1-s_{0k})$ at the pooled
share; because the pooled share depends on $\beta_k$, the simulator repeats itself from the same
seed until the two agree. \code{scripts/acceptance.py} runs the full pipeline, including the temperature loop, on
10 datasets per case.
\begin{center}
\footnotesize
\begin{tabular}{@{}lrrrr@{}}
\toprule
Case (10 datasets each) & Coverage & Error in $e^{\mu_u}$ & Bias of median & Flagged\\
\midrule
True uplift 3\%/bp & 81\% & 3.8\% & $-2.5\%$ & 2\\
True uplift 8\%/bp & 83\% & 2.2\% & $-1.3\%$ & 1\\
True uplift 1.5\%/bp & 83\% & 5.2\% & $-3.7\%$ & 3\\
3\%/bp, reactive repricing & 86\% & 4.8\% & $-4.2\%$ & 0\\
3\%/bp, segments off the additive pattern & 79\% & 3.0\% & $-2.7\%$ & 2\\
same, \code{segment_residual} on & 84\% & 3.2\% & $-2.9\%$ & 1\\
\bottomrule
\end{tabular}
\end{center}
Coverage is the share of segment-runs whose 80\% interval contains the true $u_k$ (target 70--90\%);
the error in $e^{\mu_u}$ is $|\text{median}/\text{truth}-1|$ averaged over datasets (target under
25\%); the bias is the median over segments of $u^{p50}_k/u_k-1$, averaged over datasets (target
within $\pm20\%$ under reactive repricing). All criteria are met. The temperature loop converged in
every run, in 2 to 6 iterations. Three things to note: a small downward lean of 1--4\% in every
case; coverage that ranges from 50\% to 100\% across single datasets, because errors are shared
across segments through the pooled parameters; and the unreliable-run flag firing on about 15\% of
fits, always for a handful of divergent transitions (1 to 7 in 8{,}000 draws). The simulated data follow the model's
assumptions, so this tests the estimator and the code, not the assumptions. Details are in
\code{docs/acceptance.md}.

\takeaway{The accuracy table says how wide the forecast range is for each segment, group and move
size, and whether that meets $\pm5\%$. Expect it to be met for small moves in large segments and in
groups. The model is checked by hiding past moves and predicting them, by forecasting quiet
periods, by inventing fake moves, and by rerunning under other assumptions. The first check needs
no model at all: look at what applications did after past moves.}

%=====================================================================
\section{Assumptions and limitations}\label{sec:limits}

\begin{enumerate}
\item \textbf{Where identification comes from.} Segment-specific moves in the gap, ours and
competitors', identify $\beta$ against the segment trends. Bank-wide moves identify it against the
bank-wide trend, and only because that trend is assumed smooth (\S\ref{sec:rw}). The sensitivity
run with week dummies shows how much rests on the second source.
\item \textbf{The prior.} With about ten moves the prior matters wherever \code{prior_share} is
high. Its median is a judgment (\S\ref{sec:prior}).
\item \textbf{Temperature.} $T$ is solved, not estimated, and its uncertainty is not in the
intervals (\S\ref{sec:Tloop}).
\item \textbf{Unseen competitor discounts} attenuate $\beta$ (\S\ref{sec:x}); \code{adj_bps} is the
only defence.
\item \textbf{Turnaround and capacity.} Turnaround is inside $\ell$. If a price cut causes a surge
that slows turnaround, $\beta$ is the net effect of both: right for forecasting if capacity behaves
the same way next time, too low if capacity has since improved.
\item \textbf{Announcement effects.} Rates enter at their effective date. If borrowers shift timing
around an announcement, the weeks either side move in opposite directions, and the lag shape absorbs
some of it.
\item \textbf{Fixed judgments.} $\kappa$, the priors, one half-life for all components and segments,
forward-filling competitor rates, the trailing 13-week market volume and conversion in forecasts.
\item \textbf{Model form.} Binary logit against a soft-min benchmark with constant $\beta_k$, no
outside good (our price does not change the market total), IIA among lenders, no spillovers between
segments, no competitor response to our moves, and no segment-level departure from the additive
pooling pattern unless \code{segment_residual} is switched on.
\end{enumerate}

\takeaway{The model learns from every change in our price relative to the market, and says how sure
it is. It is weakest where it leans on judgment: the starting range for sensitivity, the smoothness
of the bank-wide trend, and the competitor discounts it cannot see. The outputs and the sensitivity
runs show how much each of those matters.}

\end{document}
